The first workers are holding the wage they were promised. Nothing has been taken from their hands. They have worked through the heat, the day is over, and the owner has kept his word. Yet they cannot enjoy what they have received. They have seen what someone else received, and now their own gift seems smaller.

That is a painful kind of poverty: to have something good and lose the joy of it because someone else has it too. We can recognize the temptation without ever standing in a vineyard. Someone is welcomed warmly. Someone begins again after years away. Someone receives help we wish had come to us. Their good can begin to feel like an accusation against our own life. Why so much for them? What about all the years I have given?

Jesus tells this parable just after Peter has asked what those who left everything to follow him will receive. Jesus does promise them a reward. Their faithfulness matters. But he also tells them that the first will be last and the last first. Then comes this vineyard, where the owner keeps going out to call more workers. The promise to the first never becomes a promise that others must receive less.

At evening, the workers who came last are paid first. Each receives the full day's wage. Those hired early see it and expect more. When they receive their agreed payment, they complain about the equality. The owner has made these latecomers equal to people who bore the day's burden. That is the wound the story exposes. They measure his goodness by whether it leaves them ahead of someone else.

Saint Augustine asks us to look closely at that wage. He understands the coin as eternal life. One person's eternal life cannot last longer than another's. Augustine does not deny different glory among the saints. He says that the endless life God gives is truly endless for each. The latecomer does not receive a smaller eternity, and the faithful worker does not lose any of it when another enters.

That changes the way we hear the story. Life with God is not a loaf that becomes smaller every time another person takes a piece. Another person's salvation cannot diminish the life God promises you. The question is whether we want that shared life, or whether we want our own goodness to be proved by someone else's exclusion. The gift is full. It is our desire that needs room to grow.

Isaiah brings us to that same difficulty from another direction. The prophet calls the wicked person to abandon an old way of living and an old way of thinking, and to return to the Lord for mercy. Then he speaks of God's thoughts and ways rising above ours. Those words follow a promise of abundant forgiveness. God's height is also the height of a mercy greater than the measure we bring to it.

We may readily want that mercy for ourselves. It becomes harder when it reaches someone we resent. Yet today's psalm praises the Lord's goodness toward all his works. If we praise that goodness, we are asking to love a God whose generosity reaches beyond the people we would have chosen. The neighbor who returns to him is not taking our place. The mercy that receives the neighbor is the mercy on which we depend.

There is hope here for someone who fears it is too late to begin again. The late workers really are called, and they really do enter. Their late beginning does not make the owner's goodness small. Bring the years you cannot recover to the Lord who calls you now. Isaiah's summons is to return. The past need not become a reason to refuse the good that is still being offered.

But Augustine also asks a sharp question of the person planning to wait. The owner promised a wage. Who promised another hour? The workers called earlier did not tell him they would come when there was less work left. They came. Mercy gives courage to begin; it does not give us ownership of tomorrow. If there is a wrong to turn from, or a good we keep postponing, the call reaches us in the day we actually have.

And what of the person who has labored faithfully for years? That labor has not been wasted. Augustine pictures God himself cultivating us. His word breaks open the heart, pulls out what chokes it, and plants what can bear fruit. God is not becoming richer through our worship. We are being made fruitful through his care. Years of fidelity are years in which love can grow. We need not wish them away because another person begins later.

Paul shows us what that fruit looks like. He longs to be with Christ, yet he sees that remaining alive means useful work for others. Their need matters to him. He can desire Christ deeply and still give himself to the labor before him. His love of God takes the form of service to people. The early workers watch another person's portion. Paul asks what his own life can give.

That is the love named in the Collect, the opening prayer of this Mass: love of God and love of our neighbor. The prayer asks for obedience that leads toward life with God. Paul's willingness to remain helps us understand such obedience. It is more than holding a place in line. It is allowing another person's good to become something we actively desire and serve.

This change reaches deeper than forcing a pleasant expression onto resentment. A heart must be opened. The Gospel acclamation draws on the story of Lydia, whose heart the Lord opened to attend to Paul's words. God can reach the place where we have become closed. His word can uncover a complaint we have been repeating for years. Hearing him means letting that word address us, including the part of us that would prefer to keep comparing.

The Eucharist carries us further into this generosity. We receive the Christ who gave himself for us. Augustine speaks of his Body and Blood at the Lord's table and draws a demanding consequence: those who receive Christ's self-giving are called to give themselves for their brothers and sisters. The gift nourishes a life that can become generous. We receive from the one who loved us at the cost of his own life.

There is work here for this week. Think of one person's good that you have found difficult to welcome. Begin by thanking God for that good, even if your feelings take time to follow. Then make one quiet choice that helps the person flourish: offer needed help, speak a truthful word of encouragement, or give time without keeping an account of who deserves more. Let gratitude become something your hands can do.

The prayer after Communion asks for continuing help so that redemption bears fruit in our conduct. We still need that help when worship ends and an old resentment returns. The Lord's gift does not leave us alone with an impossible command. The giver continues to cultivate the receiver. What we receive from Christ is to become a way of living with one another.

Saint John Chrysostom says that the kingdom has no place for envy: those who give themselves for sinners now will rejoice in their salvation there. That is the joy toward which Christ calls us. Return to the workers holding their wage. Imagine receiving the promised good and being free to rejoice that another has received it too. No gift has been taken from your hands. There is simply another person with whom to share the joy. That is a freedom worth beginning today.
